% ============================================================
% 実験の目的（手引書のはじめ，注意事項の前に入れる）
% 参考にした資料:
%   基礎ゼミ実験配布テキスト.pdf「1 実験の目的」
%   レポートの書き方.pdf「1 実験の目的」
% ============================================================

\introtitle{実験の目的}

本科目では，テスターやオシロスコープを使った基本的な測定から，ダイオードやトランジスタを用いた電子回路，電圧波の観測まで，電気電子工学のさまざまな現象を実際に体験します。その中には，これまでに講義で学んだ内容だけでなく，ダイオードやトランジスタ，電圧波のように，これから講義で詳しく学ぶ内容も含まれています。講義で学んだことは，実験で確かめることで本当に納得できることが多くあります。また，これから学ぶ内容を先に実験で体験しておくことで，講義で学ぶときに具体的なイメージをもって理解できるようになります。本科目では，次のことを目的とします。

\begin{enumerate}[label=\textbf{\arabic*.}, leftmargin=*, labelsep=0.8em, itemsep=0.4\baselineskip]
  \item 講義で学んだ電気電子工学の基礎的な原理や法則を実験によって確かめ，理解をより深く確実なものにします。
  \item ダイオードやトランジスタ，電圧波など，今後の講義で学ぶ内容を実際に体験し，電気電子回路の動作について具体的なイメージを養います。
  \item テスターやオシロスコープなどの測定器の特性と使い方を知り，電圧や電流などの電気的な量を測定する基本的な技術を身につけます。
  \item 電圧・電流・抵抗・周波数など，電気電子工学で扱う基本的な量について，その大きさの感覚（数量的な概念）を身につけます。
\end{enumerate}

また，実験を通して，技術者として必要な次のような力も養います。

\begin{itemize}
  \item 実際の現象をよく観察し，得られた結果を客観的に考察する力
  \item 限られた時間と与えられた機器を有効に使って実験を進める計画性
  \item グループで協力し，それぞれが役割を果たしながら実験を進める態度
  \item 実験の内容と結果をレポートにまとめ，筋道立てて伝える力
\end{itemize}
